\documentclass{jsarticle}
\usepackage[dvipdfmx]{graphicx}
\usepackage{amsmath}
\usepackage{amssymb}
\begin{document}
\title{Explain in Global defferential equation and \\
Global integrate equation. \\
Varintegrate equation, and horizen cut of equations.}
\author{Masaaki Yamaguchi}
\date{}
\maketitle


\newcommand{\variint}{%
	\begin{picture}(15,30)
		\put(0,0){
			$ \iint $}
		\put(0,5){\line(15,0){15}}
	\end{picture} %
}

\newcommand{\variiint}{%
	\begin{picture}(20,30)
		\put(0,0){
			$ \iiint $}
		\put(0,5){\line(15,0){20}}
	\end{picture} %
}



 $$ {\partial \over \partial f}F(x) = \int \int \text{cohom}D_k(x)[I_m]  $$
 Cohomology element is constructed with isotopy of D-brane,
 and this isotopy of symmetry structure is sheaf of manifold,
 more over this brane of element is double integrate of projection.
 $$ {\partial \over \partial f}F = {{}^t \variint} \text{cohom}D_k(x)^{\ll p} $$
 And global partial defferential equation is integrate of cut in cohomolgy,
 and this step of defferential D-brane of sheap value is varint of equals,
 inverse of horizen of cute of section value is reverse of integrate of 
 operators. Three dimension of varint of manifolds in operator of sheaps,
 and this step of times of value is equal with D-brane of equations.
 Global differential equation and integrate equation is operator of 
 global scales of horizen cut and add of cut equation are routs.

 $$ \oint r dx_m = \mathcal{O}(x,y) $$

 $$ {{}^t \variiint}_{{D(\chi,x)}} \text{Hom}[D^2 \psi]^{\ll p} \cong 
 \text{vol}\left({V \over S}\right) $$

 $$ {\partial \over \partial f}F(x) = {{}^t \variint} 
 \text{cohom}D_k(x)^{\ll p} $$

 $$ {d \over df}F = (F)^{f^{'}}, \int F dx_m = (F)^{f}  $$
 $$ {d \over df}\int \int{1 \over (x \log x)^2}dx_m =
 \left(\int \int{1 \over (x \log x)^2}dx_m \right)^{f^{'}}  $$
 $$ f^{'} = (2x (\log x + \log (x + 1))) $$
  $$  \int \int F dx_m = \left({1 \over (x \log x)^2}\right)^{
	  (\int 2x (\log x + \log (x + 1))dx)} $$ 
   $$= e^{-f} $$
  $$ {d \over df}F = \left({1 \over (x \log x)^2}\right)^{
	  (2x (\log x + \log (x + 1))) } $$  
  $$= e^{f} $$
  $$ \log (x \log x) \ge 2(\sqrt{y \log y})^{1 \over 2} $$
  大域的積分方程式自体に多重積分の回数が決まっている。
  大域的微分方程式に、XのX乗のエントロピー不変量の微分量が計算に
  関わっている。外微分をこのエントロピー式に入れる。



\end{document}
